\chapter{Related Work}
\label{sec:related_work}

This section reviews and compares the literature across relevant domains to ground our methodology and findings.

\section{Tool Use and Tool Synthesis}

The ability of language models to call external functions has been studied extensively. Early work established that models can be prompted or fine-tuned to select tools and format arguments from natural-language instructions~\citep{schick2023toolformer,qin2023toolbench}. More recent evaluation suites have scaled this substantially: MCP-Universe~\citep{mcpuniverse2025} evaluates agents against eleven real-world MCP servers spanning six domains; MCP-Bench~\citep{mcpbench2025} scales this to 250 tools across twenty APIs and twenty evaluated models; and MCPAgentBench~\citep{mcpagentbench2024} benchmarks agent efficiency in tool selection and invocation within the
Model Context Protocol (MCP) ecosystem. Common to all of these
is the assumption that tools already exist: the benchmark provides a pre-built function library and evaluates how well the agent navigates it. Our work removes this assumption.

The closest prior work to ours is the \emph{LLM as Tool Maker} (LATM) framework of \citet{cai2023latm}, in which a powerful model generates reusable Python tool implementations that a lighter ``tool user'' model subsequently invokes. The key insight is a separation of concerns: tool creation is expensive (large model, done once per task class) while tool use is cheap (smaller model, done many times). Our work extends this separation to the more demanding setting of REST API integration, where generated tools must satisfy strict interface contracts enforced by live
external services. LATM evaluates generated tools through auto-generated unit tests on a small validation sample; we evaluate through end-to-end task execution against real sandboxed APIs, which catches failures that static analysis cannot --- such as HTTP parameter mismatches or authentication schema errors.

Concurrent work has also begun to close this gap. ToolFactory~\citep{toolfactory2025} proposes a pipeline that parses heterogeneous, unstructured REST API documentation and generates callable tool stubs --- explicitly targeting the
setting where formal OpenAPI schemas are unavailable. Doc2Agent~\citep{doc2agent2025} scales a similar idea, demonstrating that LLMs can produce usable tool implementations from unstructured documentation across dozens of APIs --- the closest parallel to our synthesis pipeline. Where these works
focus primarily on synthesis coverage, we additionally evaluate the synthesized servers end-to-end through agent task execution, measure documentation grounding, and systematically probe the effect of parametric knowledge through controlled ablations.

Critically, both ToolFactory and Doc2Agent generate individual tool stubs or callable functions; neither produces a complete, deployable server exposing a typed discovery protocol that an agent can invoke at runtime. Our synthesis target is a full MCP server --- including authentication, multi-tool coordination, and protocol compliance --- which represents a more demanding and practically relevant integration target.

A complementary challenge arises when APIs evolve after the synthesis stage. Recent work on evolving tool use~\citep{evolvingtoolsiclr2025} employs Monte Carlo Tree Search to adapt tool invocations when API signatures change at inference time, framing API evolution as a planning problem. This is
distinct from our setting, where we synthesize the server once per API version, but it points to the broader challenge of keeping tool infrastructure current --- the problem that motivates automated synthesis in the first place.

Related work has also explored automated client generation from formal API specifications such as OpenAPI schemas, but these approaches require structured machine-readable inputs. In practice, many APIs provide only unstructured documentation, or provide schemas that are incomplete or misaligned with actual server behaviour. Our benchmark specifically targets this less
structured setting, using scraped Markdown documentation as the sole synthesis input, which is closer to the conditions a practitioner faces.

Agent benchmarks such as WebArena~\citep{webarena2023}, OSWorld~\citep{osworld2024}, and GAIA~\citep{gaia2023} evaluate agents against fixed tool environments, taking the tool layer as given. A related line of work examines agents that generate and execute code to solve tasks~\citep{codeact2024}, but the generated code is typically a self-contained script rather than a persistent, reusable interface. We study a complementary question: can an agent produce tool infrastructure that other agents can then use?

\section{Code Generation}

Classical code-generation benchmarks such as HumanEval~\citep{chen2021codex} and SWE-bench~\citep{jimenez2023swebench} evaluate against unit tests or static oracles. These settings share a structural limitation for our purposes: they do not run generated code against a live external service, so a synthesized API wrapper that compiles and passes mock tests may silently fail when the real API rejects its HTTP requests.

Recent work has begun to study live API execution directly.
WAPIIBench~\citep{wapiibench2025} catalogues the error modes of LLM-generated API integration code against real services, identifying categories such as authentication failures, incorrect parameter serialisation, and endpoint version mismatches. This taxonomy aligns closely with the failure modes we
observe in our synthesis evaluation. RAG-augmented API code
generation~\citep{ragapidocs2025} shows that retrieving relevant documentation snippets at generation time reduces hallucinated parameters and improves execution success, suggesting that context grounding is a tractable intervention even for models with strong parametric priors. When LLMs Lag Behind~\citep{lagging2025} focuses specifically on faithfulness failures arising from training-cutoff drift --- cases where the model generates syntactically correct code that silently targets a deprecated API version, producing errors that neither compile-time nor unit-test evaluation would catch. Taken together,
these works confirm that live API execution is a qualitatively different evaluation regime from static code correctness, and that documentation faithfulness is a first-class concern rather than a secondary one.

\section{In-Context Learning and Context Faithfulness}

In-context learning (ICL), the ability of LLMs to adapt to new inputs from examples or instructions in the prompt, was systematically characterised by \citet{brown2020gpt3}. A common assumption in the software engineering community is that supplying live API documentation in the context window grounds the model's output, ensuring that generated code reflects the current API specification rather than stale parametric memory. Subsequent work has expanded on this assumption in several ways.

\citet{liu2023lostmiddle} demonstrated that recall follows a U-shaped curve across context position: performance is highest for information near the beginning or end of the context and lowest for information placed in the middle. For tool synthesis over long documentation files, this implies that endpoints described far from the beginning or end of a document may be systematically under-represented in the generated server --- a structural bias independent of parametric knowledge effects. More recent work on disentangling recall and reasoning in long contexts~\citep{disentanglerecall2025} distinguishes between a model's ability to retrieve information from the context window and its ability to reason over that information; the two can fail independently. In our setting, synthesis failures attributable to missed documentation may reflect retrieval limitations rather than reasoning failures.

More directly relevant is work on how models handle context that explicitly contradicts their parametric priors. \citet{xie2023adaptive} showed that models can be steered by stated premises even when those premises are false, while other studies have found the opposite --- models ignoring contradictory context when their parametric prior is strong. Our mutation design sits at this intersection: by introducing deliberate conflicts between in-context documentation and training-time parameter names, we measure which signal dominates, and for which APIs, providing a controlled analogue to the natural conflicts that arise when API documentation evolves after a model's training cutoff.

A recent study on task-dependent conflict resolution~\citep{taskmatters2025} finds that the degree to which a model defers to context versus its parametric prior depends not only on entity frequency but also on the nature of the downstream task: models are more likely to update on context when the task is procedural rather than factual. This is directly relevant to tool synthesis, where the model must not merely recall a parameter name but correctly propagate it through the generated code --- a procedural constraint that might be expected to promote context adherence, yet our results show this is often not the case for well-known APIs.

\section{Parametric Knowledge and Knowledge Conflicts}

A growing body of work examines how language models resolve conflicts between information encoded in their weights during pre-training and information provided in context at inference time. Studies have found that models exhibit a strong prior towards their training distribution, sometimes ignoring counter-factual evidence in the context~\citep{longpre2021entity}, and that this tendency grows stronger as the amount of training exposure to
a topic increases~\citep{mallen2023knowns}.

\citet{mallen2023knowns} provide perhaps the most directly
relevant finding: the probability that a model correctly answers a factual question from parametric memory scales with the log-popularity of the corresponding entity (measured via Wikipedia page views as a proxy for pre-training frequency). Rare entities are reliably retrieved from context; common entities are retrieved from weights even when context is available.
For our setting, this predicts that no-docs synthesis quality should be highest for APIs heavily represented in pre-training
corpora, which motivates stratifying APIs by familiarity in our benchmark design.
\citet{longpre2021entity} studied entity-based knowledge
conflicts in question answering and found that models with stronger parametric priors are more likely to ignore contradictory context --- a pattern directly relevant to the counterfactual design we adopt.

This phenomenon has also been studied in question answering settings where retrieved passages contradict model priors. \citet{neeman2022disentqa} introduce a counterfactual data augmentation approach that trains models to maintain separate predictions for parametric and contextual answers simultaneously, improving robustness to knowledge conflicts. The fact that a specialised training objective is needed to achieve this separation underscores how strongly parametric priors resist contextual override by default.

More recent work has sharpened the understanding of when parametric knowledge yields to context. JuICE~\citep{juice2025} introduces a controlled training objective that suppresses parametric interference during context-grounded generation, and demonstrates consistent accuracy gains on conflict-heavy QA benchmarks. Separately, studies on calibration in retrieval-augmented generation (RAG) pipelines~\citep{ragknowswrong2025} show that models often
fail to signal uncertainty when their parametric prior conflicts with retrieved evidence, producing confident but incorrect outputs without any warning.

In the tool synthesis setting, parametric gravity creates a measurable tension: popular APIs such as GitHub and Stripe are extensively represented in training corpora, so a model may produce a correct wrapper from memory alone --- or may resist updating that wrapper when documentation specifies something different. We build on this line of work to study how this tension plays out in a live API execution setting.

\section{Mechanistic Interpretability and Probing}

The behavioural evidence from adherence rates raises a mechanistic question: at which point in the forward pass does parametric memory override in-context information? Linear probing provides a systematic methodology for investigating this. \citet{alain2017understanding} demonstrated
that linear classifiers trained on intermediate layer representations reveal when and where concepts become linearly decodable across the network's depth, treating each layer's activation as a probe target. Applied to transformer
language models, this technique has been used to track the emergence of syntactic and semantic structure across layers~\citep{tenney2019bert}.

The logit lens~\citep{nostalgebraist2020logitlens} extends this idea directly: by projecting intermediate residual stream states through the model's unembedding matrix at each layer, it is possible to read off the model's evolving prediction, revealing how token probabilities develop from early to late layers. \citet{belrose2023tunedlens} refined this into the tuned lens, training a learned affine transformation per layer to improve the quality of intermediate predictions. These tools make it possible to detect at which layer a model's prediction transitions from an in-context representation to a parametric one.

\citet{meng2022rome} localised factual recall in GPT-style models to specific middle-layer MLP modules, finding that factual associations are stored as a form of key-value memory in feed-forward layers. This provides a candidate mechanistic account of parametric gravity: in-context representations of mutated parameter names may compete with these stored associations at the feed-forward layer where the expected parameter name is ``looked up.''

More recent work has refined this picture further. Factual recall via associative memories~\citep{factualrecall2025} models factual associations in transformers as Hopfield-style attractor networks encoded in feed-forward MLP weights, providing a principled account of why well-known entities are resilient to context
override. \citet{posttrainingreshapes2025} demonstrate that instruction fine-tuning and reinforcement learning from human feedback (RLHF) alter the geometry of the residual stream in ways that affect how strongly parametric
associations resist contextual updates --- a finding relevant to our observation that different fine-tuned models (GPT-5.2, Sonnet~4.6) show markedly different adherence rates despite comparable parametric exposure to the GitHub API.

Our probing analysis uses a log-linear classifier on the
residual stream at each layer to detect the point at which representations of mutated parameter names diverge from their training-time counterparts, connecting the behavioural adherence results to this mechanistic literature.

\section{The Model Context Protocol}

MCP~(\citeauthor{anthropic2024mcp}, \citeyear{anthropic2024mcp}) is an open standard that defines how agents discover and invoke tools exposed by external servers over a stdio transport. An MCP server exposes a \texttt{list\_tools()} endpoint that returns typed tool schemas and a \texttt{call\_tool()} endpoint for invocation. The protocol has seen rapid adoption across agent frameworks, making it a natural target for automated tool synthesis. MCP-Universe~\citep{mcpuniverse2025} evaluates agents against existing MCP servers; our work instead generates those servers from scratch and asks whether the generated infrastructure is sufficient for agent task success.

\begin{table}[h]
\centering
\scriptsize
\caption{Positioning relative to prior work. \emph{Tool synth.}: executable tool creation from documentation. \emph{Full server}: complete deployable server (not individual stubs). \emph{Doc faith.}: output grounded in provided documentation. \emph{Mech.\ analysis}: internal representation probing.}
\label{tab:related}
\setlength{\tabcolsep}{0.5pt}
\begin{tabular}{lcccccc}
\toprule
& \textbf{Tool use} & \textbf{Tool synth.} & \textbf{Full server} & \textbf{Live API}
& \textbf{Doc faith.} & \textbf{Mech.\ analysis} \\
\midrule
ToolBench        & \checkmark &            &            &            &            & \\
API-Bank         & \checkmark &            &            &            &            & \\
LATM             &            & \checkmark &            &            &            & \\
HumanEval        &            & \checkmark &            &            &            & \\
SWE-bench        &            & \checkmark &            &            &            & \\
MCPAgentBench    & \checkmark &            &            &            &            & \\
MCP-Universe     & \checkmark &            &            &            &            & \\
GAIA  & \checkmark &            &            &            &            & \\
ToolFactory      &            & \checkmark &            &            & \checkmark & \\
Doc2Agent        &            & \checkmark &            & \checkmark &            & \\
WAPIIBench       & \checkmark &            &            & \checkmark & \checkmark & \\
MCP-Bench        & \checkmark &            &            & \checkmark &            & \\
\textbf{Ours}    & \checkmark & \checkmark & \checkmark & \checkmark & \checkmark & \checkmark \\
\bottomrule
\end{tabular}
\end{table}

\section{Summary}

Table~\ref{tab:related} summarises the positioning of our work relative to the benchmarks discussed above. The key distinguishing features are: we evaluate the \emph{production} of complete, deployable tool servers (not individual stubs) from documentation combined with live API execution, and we isolate the contribution of parametric knowledge through controlled ablations --- capabilities absent from all prior work.